\section{Zadanie 16}

Mysz i kot niezależnie od siebie przemieszczają się w sposób losowy po
spójnym, nieskierowanym i niedwudzielnym grafie G.
Niech \(n\) = \( |V(G)|, m = |E(G)| \). Pokaż, że \(O(m^2n)\) jest
ograniczeniem górnym na oczekiwaną liczbę kroków, po której dojdzie do
spotkania.

\textit{Wskazówka: rozważ łańcuch Markowa na parach indeksów
	\((a, b)\)}

\subsection*{Rozwiązanie}

Oznaczmy wierzchołki i krawędzie naszego grafu \(V\) i \(E\)
(czyli \(G = (V, E)\)).

Zbudujemy graf \(H = (V_H, E_H)\), gdzie
\begin{itemize}
	\item \(V_H = V \times V\) - wierzchołek nowego grafu będzie parą
	      wierzchołków w oryginalnym grafie
	\item \(E_H = \big\{\{(u, v), (w, t)\} \in \binom{V_H}{2} \colon
	      \{u, w\} \in E \land \{v, t\} \in E
	      \big\}\)
	      \\- krawędź nowego grafu będzie oznaczała, że w grafie \(G\) były takie
	      dwie krawędzie, którymi można było przejść z jednej pary wierzchołków
	      do drugiej.
\end{itemize}
Zauważmy, że spacer losowy w grafie \(H\) jest tym samym co operacja
z treści zadania.
Załóżmy, że mysz i kot byli w wierzchołkach \(a\) i \(b\), a po
ruchu są w \(c\) i \(d\). Szansa na to w oryginalnym grafie to
\(\frac{1}{d(a)}\cdot\frac{1}{d(b)}\) (z definicji spaceru losowego).
\\W grafie
H oznacza to przejście z wierzchołka \((a, b)\) do \((c, d)\),
czyli szansa była \(\frac{1}{d_H((a,b))}\).\\
\(d_H((a, b))\) jest liczbą par wierzchołków, do których można się
dostać z pary \((a, b)\), czyli dokładnie \(d(a)\cdot d(b)\).\\
Prawdopodobieństwo tych dwóch zdarzeń jest identyczne.

Łatwo zauważyć, że \(|V_H| = |V| \cdot |V| = n^2\)\\
Trudniej \(|E_H| = \frac{1}{2} \sum_{a, b \in V} d(a)d(b)
= \frac{1}{2} \sum_{a \in V} d(a) \sum_{b \in V} d(b)
= \frac{1}{2} \cdot 2m \cdot 2m = 2m^2
\)

Znajdźmy w grafie \(G\) nieparzysty cykl \(c\)
(graf nie jest dwudzielny).

Załóżmy, że mysz zaczęła w wierzchołku \(x\), a kot w \(y\).

Niech \(p_1\) oznacza spacer (ścieżkę z możliwymi powtórzeniami)
długości \(n\) zaczynający się w \(x\) i kończący się w \(c\).
Taki spacer oczywiście istnieje. W grafie spójnym w co najwyżej
\(n-1\) krokach można się dostać do dowolnego wierzchołka,
a następnie można krążyć wokół cyklu.

Analogicznie niech \(p_2\) oznacza spacer z \(x\) do dowolnego
punktu na \(c\), długości dokładnie \(n\).

Teraz kontynuujmy spacery \(p_1\) i \(p_2\) równocześnie w
przeciwnych kierunkach wokół cyklu. Po co najwyżej
\( |c| \leq n \) krokach spacery spotkają się
(bo \(c\) miał nieparzystą długość)

Znaleźliśmy dwa spacery o równej długości, nie większej niż \(2n\),
które kończą się w tym samym punkcie.
Te dwa spacery w grafie \(G\) możemy potraktować jako jeden spacer w
grafie \(H\), zaczynający się w wierzchołku \((x, y) \in V_H\), a
kończący się w jakimś \((z, z)\). Oznaczmy go \(p\).

Skorzystamy z lematu z ćwiczeń (\texttt{Lemma 7.14}):
Dla krawędzi \((u, v)\) mamy \(h_{u, v} + h_{v, u} \leq 2|E|\),
gdzie \(h_{u, v}\) to oczekiwany czas osiągnięcia wierzchołka \(v\)
zaczynając w \(u\).\\
Nam wystarczy \(h_{u, v} \leq 2|E|\).

Będziemy iterować się po krawędziach na ścieżce \(p\) w grafie \(H\).
Oczekiwany czas przejścia jedną krawędzią jest z lematu
co najwyżej \(2|E_H|\), a krawędzi jest co najwyżej \(2n\).

\[
	\ev{\texttt{liczba kroków}} \leq |p| \cdot 2|E_H|
	\leq 2n \cdot 2 \cdot 2m^2 = 8m^2n = O(m^2n)
	\quad \square
\]
